% =====================================================================
% EDS Acquisition Parameters for CALIBER - compiled ideas & rationale
% Source of truth for the PDF docs/eds_parameter_ideas.pdf
% Build: pdflatex -interaction=nonstopmode eds_parameter_ideas.tex (x2)
% To request an update: tag @claude on the PR holding this file (or any
% PR/issue) with the info to add; a new revision row is added per update.
% =====================================================================
\documentclass[11pt]{article}
\usepackage[margin=0.95in]{geometry}
\usepackage[T1]{fontenc}
\usepackage[utf8]{inputenc}
\usepackage{lmodern}
\usepackage{amsmath}
\usepackage{amssymb}
\usepackage{xcolor}
\usepackage{enumitem}
\usepackage{booktabs}
\usepackage{parskip}
\usepackage{array}
\usepackage[colorlinks=true,linkcolor=blue!55!black,urlcolor=blue!55!black,citecolor=blue!55!black]{hyperref}

% ------- attribution tags -------
\definecolor{cR}{RGB}{23,111,44}    % Ronnie - green
\definecolor{cC}{RGB}{28,80,158}    % Claude - blue
\definecolor{cRC}{RGB}{0,110,110}   % Ronnie -> Claude - teal
\definecolor{cCR}{RGB}{110,45,150}  % Claude -> Ronnie - violet
\definecolor{cMS}{RGB}{160,82,10}   % Mike Standing via Ronnie - brown
\newcommand{\tR}{\textcolor{cR}{\textbf{[Ronnie]}}}
\newcommand{\tC}{\textcolor{cC}{\textbf{[Claude]}}}
\newcommand{\tRC}{\textcolor{cRC}{\textbf{[Ronnie $\to$ Claude]}}}
\newcommand{\tCR}{\textcolor{cCR}{\textbf{[Claude $\to$ Ronnie]}}}
\newcommand{\tMS}{\textcolor{cMS}{\textbf{[Mike Standing $\to$ Ronnie]}}}

% ------- provenance link helpers -------
\newcommand{\pr}[1]{\href{https://github.com/vertical-cloud-lab/caliber/pull/#1}{PR~\##1}}
\newcommand{\prc}[2]{\href{https://github.com/vertical-cloud-lab/caliber/pull/#1\#issuecomment-#2}{PR~\##1 c-#2}}
\newcommand{\issc}[1]{\href{https://github.com/vertical-cloud-lab/caliber/issues/1\#issuecomment-#1}{issue~\#1 c-#1}}
\newcommand{\basis}[1]{\par\nopagebreak{\footnotesize\emph{Basis / provenance:} #1\par}}

\setlist[itemize]{itemsep=5pt, topsep=3pt, leftmargin=1.4em}

\title{\textbf{EDS Acquisition Parameters for CALIBER}\\[4pt]
\large Compiled ideas, attributions, and the reasoning behind them\\
\normalsize Beam energy \,\textperiodcentered\, Acquisition time \,\textperiodcentered\, Beam current --- AlSi10Mg first campaign}
\author{Compiled by Claude from \pr{11}, \pr{13}, \pr{15}, and the
\href{https://github.com/vertical-cloud-lab/caliber/issues/1}{issue~\#1} discussion\\
\texttt{vertical-cloud-lab/caliber}}
\date{September 16, 2026 --- rev.\ 1.0}

\begin{document}
\maketitle

\section{How to read this document}

This document collects every idea voiced in the EDS parameter threads --- \pr{11}
(parameter 1, beam energy), \pr{13} (parameter 2, counts / acquisition time),
\pr{15} (parameter 3, beam current) --- plus ideas proposed by Claude that Ronnie
confirmed for the program. ``Ideas'' means the reasons that go into choosing the
\emph{initial} optimal value of each parameter, and the reasons that govern
\emph{optimizing it afterwards}. Each idea carries an attribution tag, the physical
reasoning or paper behind it, and a link to where it was said.

\subsection*{Attribution legend}
\begin{itemize}[itemsep=1pt]
  \item \tR{} --- idea voiced by Ronnie Guymon.
  \item \tC{} --- idea contributed by Claude.
  \item \tRC{} --- idea originated by Ronnie, then quantified, refined, or verified by Claude.
  \item \tCR{} --- proposed by Claude, then confirmed or adopted by Ronnie for the program.
  \item \tMS{} --- advice from Mike Standing (BYU SEM specialist), relayed through Ronnie's notes.
\end{itemize}

\subsection*{Living document}
\noindent\fbox{\parbox{0.965\textwidth}{\small
The LaTeX source \texttt{docs/eds\_parameter\_ideas.tex} lives next to this PDF in the
repository, so any \texttt{@claude} session can extend it and recompile. To add
material, tag \texttt{@claude} with the info and say it is for the \emph{EDS parameter
ideas PDF} --- ideally on the pull request that holds this file (updates then land on
the same branch), but a ping from any other PR or issue works too. Every update adds
a row to the revision table below.}}

\begin{center}
\begin{tabular}{lll}
\toprule
rev. & date & change \\
\midrule
1.0 & 2026-09-16 & Initial compilation from PRs \#11, \#13, \#15 and issue \#1. \\
\bottomrule
\end{tabular}
\end{center}

\section{The optimization framing these parameters live in}

\begin{itemize}
\item \tR{} (with Xavier Zaitzeff) \textbf{Objective function.} Minimize
$(\mathrm{EDX} - \mathrm{ICPMS})^2$: EDS-measured composition judged against an
ICP-MS ground truth. Search space from the planning outline: beam energy/voltage
$\{5, 10, 15, 20\}$~kV; acquisition time $[1~\mathrm{min},\, 30~\mathrm{min}]$;
beam current $[20~\mathrm{pA},\, 10~\mathrm{mA}]$ (range marked with a ``?'' ---
to be verified against the instrument, see \S\ref{sec:current}); 16 Sobol samples
for the initial design; a Gaussian-process chain (microscope parameters
$\to \mathrm{GP}_1 \to Y_1 \to \mathrm{GP}_2$ with post-processing $\to$ objective);
later, tests on other alloys.
\basis{Whiteboard outline by Ronnie and Xavier, committed as
\href{https://github.com/vertical-cloud-lab/caliber/blob/46140a4/CALIBER_Outline.jpg}{\texttt{CALIBER\_Outline.jpg}}.}

\item \tR{} (with Xavier; Gage and Dr.\ Rappleye arranging) \textbf{One-time ICP-MS
ground truth.} Run ICP-MS once on the tensile sample to obtain the ``truth'' value
that all EDS measurements are compared against --- the yardstick for judging whether
parameter optimization is actually improving accuracy.
\basis{\issc{5574532841}; Claude had earlier suggested an ICP referee measurement
(ICP-OES) when the 0.94 vs 0.5--0.7 wt\% gap appeared \tC{}. Expected truth window
for Mg: powder certificate
0.5--0.7~wt\% (the earlier 0.05--0.07 figure was a typo, corrected by Ronnie);
Edison literature task on Mg evaporation in LPBF predicts near-complete retention
under Ar, so $\approx$0.45--0.70~wt\% as-built
(\href{https://github.com/vertical-cloud-lab/caliber/blob/d6c29d1/outputs/issue-1-mg-evaporation/answer.md}{Edison artifacts}).}

\item \tCR{} \textbf{Standards-based, un-normalized quantification is the measurement
backbone.} Vendor standardless output is normalized to 100\%, which entangles every
element's error with every other's (the bogus 8.1~wt\% O row at 15~kV dragged Si
down); an un-normalized analytical total (98--102\% = healthy) is itself a QC check.
Claude argued the physics; Ronnie adopted the DTSA-II standards route and verified a
manual k-ratio pass himself (Mg $\approx$ 0.9~wt\%).
\basis{Newbury, Swyt \& Myklebust [11] on standardless risk; Newbury \& Ritchie [2]
for the k-ratio protocol; \prc{11}{5481281979} \S3; Ronnie's DTSA-II pass
\prc{11}{5373657846}.}
\end{itemize}

% =====================================================================
\section{Parameter 1 --- Beam energy (accelerating voltage, $E_0$)}

\subsection{Why this parameter dominates}

\tR{} Ronnie's synthesis (\issc{5496243634}), verbatim:

\begin{quote}\small
``The parameter that seems to have the largest impact is the beam energy, since it
effects spatial resolution, the amount of correction needed to filter through the
background noise, intensity of the x-ray spectra lines, absorption (from path length
and elemental edge absorption), peak crowding, and homogeneity.''
\end{quote}

Claude's additions to that axis list: \textbf{surface-film weighting} (the main force
pushing \emph{against} low kV), \textbf{count rate / dead-time management}, and
\textbf{standards logistics per voltage}. One consolidation \tC{}: ``edge absorption''
and ``path-length absorption'' are one correction with two knobs, not two corrections
--- attenuation is $\exp[-(\mu/\rho)\,(\rho z/\sin\psi)]$, where the absorption edges
set $\mu/\rho$ (chemistry only; kV never touches it) and kV sets the depth $\rho z$.
They multiply inside the same exponent.
\basis{\prc{11}{5481278418} (Ronnie's list); \prc{11}{5481281979} \S4 (additions).}

\subsection{Choosing the initial voltage}

\begin{itemize}
\item \tRC{} \textbf{The slant-path absorption criterion --- the core CALIBER rule.}
Ronnie's proposal: choose a voltage whose interaction volume, at its farthest point
from the surface, is shorter than the absorption path length of Si. Claude's two
refinements: (a) the photon exits along the detector direction, so the comparison is
\emph{exit path} = depth $\times \csc\psi$ = depth $\times 1.74$ at the Apreo's
$\psi = 35.1^\circ$ take-off, not the depth itself; (b) use the X-ray
\emph{production} depth (Kanaya--Okayama range with $E_0^{1.67} \to E_0^{1.67} -
E_c^{1.67}$, since electrons below the edge energy $E_c$ produce nothing), not the
full electron range. The rule:
\[ (\text{deepest production depth}) \times \csc\psi \;\le\; (\text{$1/e$ attenuation length}) \]
equivalently $f(\chi) \gtrsim 0.75$ for the worst quantified line. It pins every
photon's absorption probability below 63.2\%; the spectrum-average absorption at the
cap is only about 25\%. The slant factor is load-bearing: with it the AlSi10Mg cap is
\textbf{7.6~kV} (only 5~kV passes on the fixed menu); without it, 10.3~kV would just
admit 10~kV.
\basis{Ronnie's proposal \prc{11}{5481278418}; Claude's computation and refinement
\prc{11}{5481281979} \S4 and \prc{11}{5498155829}; Kanaya \& Okayama [6]; Philibert
$f(\chi)$ with Elam MACs [7]; independently recomputed in
\texttt{scripts/eds\_voltage\_predictions.py} (\pr{13}), reproducing all values.}

\item \tRC{} \textbf{What the rule protects: correction size, not counts.} Ronnie
asked how impactful the rule is to counts; the counterintuitive answer is that it
\emph{costs} counts --- escaped Si~K$\alpha$ per unit dose still rises with kV
($\times$3.6 / 5.6 / 6.9 at 10/15/20~kV relative to 5~kV). What it bounds is the
absorption (A) correction: every absorbed count is arithmetically re-created using
mass absorption coefficients that are themselves uncertain at the 5--25\% level, so
the absorbed fraction is the part of the answer that is \emph{computed rather than
measured}. A $\pm$20\% MAC error moves reported Si by 2.6\% at 5~kV versus 16\% at
20~kV. Counts are the resource in surplus; correction accuracy is the one that is not.
\basis{\prc{11}{5498155829}; MAC sensitivity columns in
\texttt{eds\_voltage\_predictions.csv} (\pr{13}).}

\item \tRC{} \textbf{Low-energy lines magnify absorption error.} Ronnie's hypothesis
(``the lower energy the line, the greater the A-factor will magnify error present''
--- flagged for double-checking): confirmed. Softer lines have shorter attenuation
lengths, so the same model/MAC uncertainty leverages harder. O~K$\alpha$'s cap is
\textbf{4.6~kV} --- it fails even at 5~kV --- which is the criterion independently
rediscovering Statham's advice not to quantify lines below 1~keV.
\basis{Ronnie's hypothesis, \issc{5496243634}; Statham [3]; per-line caps in
\S\ref{sec:voltagetable}.}

\item \tR{} \textbf{Overvoltage floor.} A line must be excited at all
($U_0 = E_0/E_c > 1$), and Ronnie's rule: define an absolute lower bound --- if
$U_0 \lesssim 1.5$, move to the next voltage on the menu. His worked example:
Mn~K$\alpha$ (5.898~keV) gives $U_0 = 1.695$ at 10~kV --- to be evaluated whether
that suffices. Limiting case: Fe~K$\alpha$ has $U_0 = 0.70$ at 5~kV --- physically
absent (prediction P2).
\basis{\prc{11}{5482266968}; ionization cross-section peaks near $U \approx 2$--3,
behind the classic ``$E_0 \approx 2\times$ the highest line energy'' rule of thumb
(Goldstein [1]), \tC{} in the original issue-\#1 answer.}

\item \tR{} \textbf{Line-menu constraints from the element roadmap.} Planned alloying
elements: Al, Mn, Cr, Zr, Mg, Si, Cu, Ti, Fe, Ni, Ce, Sc, Li, Er, Zn, Sn. Sticking to
the below-1-keV exclusion forces K$\alpha$ for Ti/Cr/Mn (their L$\alpha$ lines sit
under 1~keV), which demands high kV; Li~K$\alpha$ at 52~eV is unquantifiable by EDS
--- a different technique is needed if Li enters the roadmap.
\basis{\prc{11}{5481972819} (roadmap), \prc{11}{5482266968} (K-line forcing, Li).}

\item \tC{} \textbf{Trace heavy elements on L lines at 5~kV do not work.} A 3d-metal
L$\alpha$ at 5~kV runs 5--10$\times$ weaker per unit concentration with 1.5--3$\times$
more background under it: detection limit of order 0.3--1~wt\% at the current dose ---
workable for \emph{minor} constituents, not trace. Statham's quantitative version
(located by Ronnie in the assigned reading): 0.001 mass fraction via L lines at 5~kV
demands background accuracy near 1\%, versus near 10\% for K lines at conventional kV
--- ``possible but ten times harder.'' This is the concrete case for a second,
higher-kV condition once heavy alloying elements arrive.
\basis{\prc{11}{5481281979} \S5; Statham [3] p.~3 (Ronnie's find,
\prc{11}{5445481577}); Currie framework [4].}

\item \tR{} \textbf{Background mass fraction and peak crowding at low kV.} From
Statham p.~3: at low kV the bremsstrahlung background is a much larger mass fraction
of the spectrum, so background-correction errors and competing peaks hurt more. And a
0--5~keV display range crowds all peaks into a window 3--4$\times$ smaller than a
15--20~kV range, raising overlap likelihood.
\basis{\prc{11}{5445481577}; Statham [3].}

\item \tC{} \textbf{Surface-film weighting is the main counter-pressure against low
kV.} At 5~kV the Kanaya--Okayama range in Al is 0.42~$\mu$m (vs 2.6~$\mu$m at 15~kV),
so carbon contamination, native oxide, and any Mg-enriched oxide skin are
over-weighted in the sampled volume (the Aug run's C 1.99 / O 1.68~wt\% prove the
layer is there). A colloidal-silica final polish leaves embedded SiO$_2$ that reads
as extra Si + O precisely at this depth --- record the final polishing step.
Mitigation adopted \tR{} (with Gage): prepare a freshly polished sample and store it
in the vacuum chamber between polish and session.
\basis{Claude's warning quoted, with Ronnie's mitigation, in \prc{11}{5374055623};
\prc{11}{5481281979} \S4; Goldstein [1] ch.~3 for ranges.}

\item \tRC{} \textbf{Interaction volume vs.\ microstructural homogeneity.} Ronnie
(from Goldstein \S3.3.3.1: interaction-volume size is a strong function of beam
energy) explored running \emph{all} readings at one high voltage so every element and
standard shares an interaction volume. Claude's counterpoint for LPBF AlSi10Mg: the
alloy is heterogeneous at the sub-$\mu$m cellular-segregation scale, so 5~kV
(0.42~$\mu$m $\approx$ cell scale) resolves Si-rich vs Al-rich phases while 15--20~kV
averages them --- which element segregates (here Si, with Mg tracking it into the
eutectic network) determines what the averaging volume does to a map. Phase-contrast
collapse with kV is prediction P4.
\basis{\prc{11}{5441452199}; Goldstein [1] \S3.3.3.1; Aug phase split (17.3 vs
7.5~wt\% Si) in \texttt{docs/eds\_parameter\_summary.md} \S2.}

\item \tC{} \textbf{Proposed long-term menu: two points, 5 + 20~kV} (status:
proposed, pending the element-roadmap conversation). 5~kV as the light-element /
small-correction / high-resolution leg; 20~kV as the transition-metal-K$\alpha$ /
trace-detection leg, knowingly outside the absorption criterion --- the accepted
trade when the only usable line needs the energy, and exactly where measured
standards earn their keep. 15~kV underexcites Cu and Zn K ($U_0$ = 1.67 / 1.55 vs
2.23 / 2.07 at 20~kV); 10~kV adds little the endpoints do not span. Cost is small:
a three-item standards block is about 15 minutes per additional voltage, but DTSA-II
binds every standard to detector + beam energy ($\pm$1\%), so each menu voltage needs
its own bundle set --- Ronnie's related note: new elements may mean standards at each
of 5/10/15/20~kV.
\basis{\prc{11}{5481281979} \S6; Ronnie's standards-burden note \prc{11}{5481278418};
fixed-menu constraint $\{5,10,15,20\}$~kV \tR{} \prc{11}{5481700165}.}
\end{itemize}

\subsection{The criterion quantified for AlSi10Mg}
\label{sec:voltagetable}

Take-off $\psi = 35.1^\circ$ ($\csc\psi = 1.74$), $\rho = 2.67$~g/cm$^3$, Elam MACs.
Values from the \pr{13} prediction script's CSV
(\texttt{eds\_voltage\_predictions.csv}), which reproduced the \pr{11} \S4
derivation exactly.

\begin{center}
\small
\begin{tabular}{lccccccc}
\toprule
line & $E$ (keV) & atten.\ length ($\mu$m) & cap (kV) &
$f(\chi)$@5 & @10 & @15 & @20 \\
\midrule
O K$\alpha$  & 0.525 & 0.61 & 4.6  & 0.72 & 0.42 & 0.25 & 0.15 \\
Mg K$\alpha$ & 1.254 & 5.60 & 17.1 & 0.96 & 0.88 & 0.79 & 0.70 \\
Al K$\alpha$ & 1.487 & 8.43 & 21.9 & 0.98 & 0.92 & 0.85 & 0.78 \\
Si K$\alpha$ & 1.740 & 1.32 & \textbf{7.6} & 0.87 & 0.64 & 0.45 & 0.32 \\
Fe K$\alpha$ & 6.404 & 38.1 & ($U_0{=}0.70$ at 5~kV) & --- & 0.99 & 0.97 & 0.95 \\
\bottomrule
\end{tabular}
\end{center}

Si~K$\alpha$ is the binding line (the ``nasty coincidence'': its energy sits just
above the Al K absorption edge, so the Al matrix absorbs it strongly --- attenuation
length only 1.32~$\mu$m). On the fixed menu, \textbf{only 5~kV complies}; O~K$\alpha$
complies nowhere.

\subsection{Optimizing / validating the voltage afterwards}

\begin{itemize}
\item \tR{} \textbf{The falsification experiment.} ``For the first AlSi10Mg
standards-based run, I want to do EDS at 5, 10, 15, and 20~kV and see if I can
actually see the effects that theoretically should come from different voltages
\ldots\ Test this hypothesis and if CALIBER worked right.'' The first standards-based
test of a CALIBER recommendation --- the recommendation must be falsifiable, not just
plausible.
\basis{\issc{5496243634}; executed as the pre-registered plan in
\href{https://github.com/vertical-cloud-lab/caliber/blob/a68e0c6/docs/eds_voltage_series_plan.md}{\texttt{docs/eds\_voltage\_series\_plan.md}} (\pr{13}).}

\item \tC{} \textbf{Pre-registered predictions P1--P7} (stated before data exist):
P1 --- kV-invariance is the core test: true composition cannot depend on kV, so any
drift in standards-based wt\% across 5$\to$20~kV is systematic error; the $\pm$20\%
MAC band widens $\pm$2\% $\to$ $\pm$16\% relative. P2 --- Fe~K$\alpha$ physically
absent at 5~kV. P3 --- apparent surface C/O diluted about 10$\times$ by 20~kV
(production depth 0.34 $\to$ 4.1~$\mu$m). P4 --- phase contrast collapses with kV.
P5 --- dead time climbs with kV at fixed current; each kV gets its own current for
the target band. P6 --- Duane--Hunt endpoint equals set kV (charging check). P7 ---
Mg/Al net-intensity ratio falls $\times$0.86 from 5$\to$15~kV. The graded stress
test: at 10~kV only Si fails (criterion ratio 1.6); at 15~kV Si fails at 3.3$\times$
and Mg turns marginal; at 20~kV Si is at 5.5$\times$ and Mg fails too --- if error
grows in that order, the criterion captures the right physics.
\basis{Plan \S2--3 (\pr{13}); prediction bands from
\texttt{eds\_voltage\_predictions.csv}.}

\item \tC{} \textbf{Pre-registered decision rules H1/H2/H3} --- the update rule is
written down before the data arrive: H1 criterion validated $\to$ CALIBER encodes
$f(\chi) \ge 0.75$ as-is; H2 too conservative (results kV-invariant even at
$f(\chi) = 0.32$ because $\varphi(\rho z)$ absorbs the physics) $\to$ relax the
floor and higher kV wins on rate, Fe access, surface robustness; H3 wrong axis
(5~kV \emph{least} accurate because surface films bias the shallow volume more than
absorption uncertainty biases the deep one) $\to$ add a surface-film weighting term.
Any of the three outcomes is useful.
\basis{Plan \S5 (\pr{13}).}

\item \tRC{} \textbf{The physics consistency check as a data-quality tool.} For the
same sample, net-intensity ratios must track generation + absorption physics when
only kV changes. Applied to Apr (15~kV) vs Aug (5~kV): Si/Al $\times$0.71 measured
vs $\times$0.66 predicted (tracks); Mg/Al $\times$1.52 vs $\times$0.86 (1.8$\times$
excess --- fails); O/Al $\times$1.33 vs $\times$0.29 (fails). Ronnie's O-normalization
hypothesis --- the bogus O row soaking up mass in the 100\%-normalization --- was
quantified as about half of the apparent 15~kV Si ``improvement''; the rest fails the
consistency check. One element moving nearer nominal while another moves 2$\times$
farther is the signature of compensating errors, not better physics. The later
dead-time reinterpretation (both sessions ran the processor $\approx$97\% idle, so
the pile-up explanation weakened) leaves the $\times$1.52 anomaly standing for P7 to
adjudicate.
\basis{\prc{11}{5481281979} \S2--3 (Ronnie's hypothesis in \prc{11}{5481278418});
dead-time reinterpretation \prc{13}{5574393906}.}
\end{itemize}

% =====================================================================
\section{Parameter 2 --- Acquisition time (live time / counts)}

\subsection{Choosing the initial acquisition time}

\begin{itemize}
\item \tR{} \textbf{Count more instead of raising voltage.} ``Instead of increasing
voltage to get more counts, we just count more.'' The rate lost by honoring the 5~kV
criterion is bought back with time and current --- not by re-raising kV.
\basis{\prc{13}{5500446306}.}

\item \tRC{} \textbf{Peak vs.\ background growth --- the corrected version.} Ronnie
and Mike both reasoned that the Mg peak should grow faster than the background,
making it more prominent. Claude's correction: peak and bremsstrahlung background are
\emph{both} proportional to delivered dose, so the peak-to-background \emph{ratio} is
fixed by kV and matrix. What longer counting improves is the peak against the
background's statistical \emph{scatter} ($\sqrt{B}$): significance and detection
limit improve as $1/\sqrt{t}$. Mg ``pops out'' in the sense of significance, not as
a cleaner ratio. Likewise the Al~K$\alpha$ tail under Mg~K$\alpha$ is a fixed
fraction of Al counts --- tripling live time triples both; the tail-model bias is
systematic and does not average away. That is what standards fix, not time.
\basis{\prc{13}{5500449993}; Mike's version in his notes, \issc{5513720218};
Ziebold [5].}

\item \tMS{} + \tC{} \textbf{The 3$\sigma$ detection criterion.} Mike: a trace
element needs to be 3 standard deviations above the background to be trusted.
Claude's formalization (Currie): decision limit $L_C = 2.33\sqrt{B}$, detection limit
$L_D = 2.71 + 4.65\sqrt{B}$ ($\approx 3.29\sqrt{B}$), with the common EDS shorthand
net $> 3\sqrt{2B}$; reliable \emph{quantification} conventionally wants
$\approx$10$\sigma$. Status for this program: the Aug 5~kV spectrum already has
Mg~K$\alpha$ at 28.2$\sigma$ with $L_D \approx 0.06$~wt\% --- detection was never
the problem; precision and accuracy are.
\basis{Mike's notes via Ronnie, \issc{5513720218} and \prc{13}{5516384978};
Currie [4]; fit numbers from \texttt{scripts/eds\_kfactor\_quant.py} (\pr{11}).}

\item \tRC{} \textbf{Count targets from the textbook, translated to this alloy.}
Ronnie located the Goldstein passage: $>10^6$ total counts enable
standards-based accuracy for major/minor constituents \emph{even under severe peak
interference}; $>10^7$ for clean 200-ppm trace work; $>10^8$ for trace under
interference. Claude's interpretation: ``interference'' means spectral peak overlap
(not VPSEM gas) --- and Mg~K$\alpha$ under the Al~K$\alpha$ tail at a 100--200:1
concentration ratio is the textbook's exact scenario, so the Mg quant lives in the
$>10^6$ tier. Unit conversion: $\ge 10^4$ net Mg~K$\alpha$ counts ($\approx$1\%
counting precision) $\Leftrightarrow$ $\approx$1.5M total counts at 5~kV
$\Leftrightarrow$ $\approx$12.5~min live at the Aug rate (2.0~kcps) or
$\approx$75~s at the rate a proper beam current delivers. Neither session reached
$10^7$ (best: 4.4M) \tR{}.
\basis{Ronnie's textbook find \prc{13}{5574391518}; Goldstein [1]; translation
\prc{13}{5574393906}.}

\item \tR{} \textbf{The cost axis.} Do not count excessively --- microscope time is
the budget; find where more data stops helping.
\basis{\prc{13}{5516384978}.}

\item \tC{} \textbf{Diminishing-returns math and a stop rule.} Benefit scales as
$\sqrt{t}$, cost as $t$: halving the statistical error costs 4$\times$ the time.
Stop when Poisson error is small next to the systematic floor (standards, matrix
correction, tail model --- realistically 1--2\% relative). Operational test: when
spot-to-spot scatter exceeds the fit-reported Poisson $\sigma$, more counting is
pure cost. Pre-registrable CALIBER rule: $\sigma_{\mathrm{stat}} \le
\tfrac{1}{3}\sigma_{\mathrm{syst}}$.
\basis{\prc{13}{5500449993}; Ziebold [5].}

\item \tC{} \textbf{Live time, not real time, is the statistics clock.} Acquisitions
are specified in live time (dead-time-corrected); real time = live time /
$(1 - \mathrm{DT})$. Counts = (input rate) $\times$ (live time) at any amp time.
\basis{\texttt{docs/eds\_daq.md} (\pr{6}, via \pr{15}); \prc{13}{5574393906}.}

\item \tRC{} \textbf{The Apr-vs-Aug lesson: rate, not time, drove the count gap.}
Ronnie compiled the sessions (Apr 15~kV: 4{,}435{,}734 counts / 460.8~s live /
1.92~$\mu$s amp; Aug 5~kV: 652{,}397 / 327.7~s / 7.68~$\mu$s) and asked why 130
extra seconds gave 3.8M extra counts. It did not --- the rates differ 4.8$\times$
(9{,}626 vs 1{,}991 cps, set by the kV: overvoltage and continuum range), times
1.41$\times$ live time = the 6.8$\times$ count ratio exactly. Plan time budgets
per kV from rate, not habit.
\basis{\prc{13}{5547583835} (Ronnie's data); \prc{13}{5574393906} (arithmetic).}
\end{itemize}

\subsection{Optimizing acquisition time afterwards}

\begin{itemize}
\item \tR{} (agreed with Paul, Mike, and Felipe; executing with Xavier)
\textbf{1M-count baseline and a doubling ladder.} Baseline of $10^6$ total counts;
collect spectra at 1M, then increase by factors of 2 (1M $\to$ 2M $\to$ 4M \ldots)
to locate the point of diminishing returns, judged against a standard --- and
ultimately against the ICP-MS truth to measure how much accuracy actually increases
with counts. First session (2026-09-17): 1/2/4M at 5~kV on the old sample, also
gauging how much carbon deposition grows as acquisition time increases.
\basis{\prc{13}{5609952233}, \prc{13}{5691161176}; converges with the Goldstein
$>10^6$ tier above.}

\item \tCR{} \textbf{Replicates instead of one long acquisition.} $5 \times 60$~s
summed is statistically identical to $1 \times 300$~s (Poisson counts add) and
retires the four real downsides of long runs: systematic-error floor, probe-current
drift, spatial drift ($\approx$10~nm/min walks a 0.3~$\mu$m probe a probe-width in
30~min), and calibration drift. Ronnie's adopted version adds a twist: multiple
shorter scans on \emph{different areas close to one another}, compiled together ---
which also spreads the carbon dose, especially valuable at low kV. DTSA-II automates
the compilation: the bundle wizard sums replicates with correct dose bookkeeping
(live times add; current becomes the live-time-weighted average) plus per-spectrum
outlier scores and a Duane--Hunt charging check; the quant wizard returns mean
$\pm$ SD across unknowns.
\basis{Claude's replicate case \prc{11}{5517130995} and \prc{11}{5515608000};
Ronnie's adoption \prc{13}{5534896987}; DTSA-II mechanics in
\texttt{docs/eds\_parameter\_summary.md} \S4.}

\item \tMS{} \textbf{Contamination management makes longer total acquisition viable.}
Longer acquisition helps statistics; the carbon-deposition cost is contained with the
facility's plasma cleaner and bake-out oven --- bake samples 2--3~h and load them into
the microscope hot. The cleaning only helps so much: cleanliness discipline starts
right after polishing.
\basis{Mike's notes via Ronnie, \issc{5513720218}.}

\item \tC{} \textbf{Measure the carbon budget instead of assuming it.} Consecutive
short spectra on one spot; watch the C~K$\alpha$ slope and Mg/Al drift. The
dose-per-spot that keeps that drift below the counting error is the third
pre-registrable CALIBER number. (C~K$\alpha$ is also where short amp times bite
hardest --- keep 7.68~$\mu$s for carbon-budget spectra.)
\basis{\prc{13}{5500449993}; amp-time caveat \prc{13}{5593623759}.}

\item \tR{} \textbf{TEAM operation.} Spectrum Only $\to$ Collect Spectrum appears to
support a counts-based stop (possibly capped near 1.1M --- Mike to confirm); the
EDAX engine's ``Limit By'' options are live time / clock time / ROI counts \tC{},
so an ROI over Mg~K$\alpha$ with a $10^4$ stop would automate the net-count target
directly. Map duration is pixels $\times$ dwell $\times$ frames, not a typed time.
\basis{\prc{13}{5593650379} (Ronnie); \prc{13}{5574393906} \S4 (TEAM references).}
\end{itemize}

% =====================================================================
\section{Parameter 3 --- Beam current (and its partner knob, amp time)}
\label{sec:current}

\subsection{Choosing the initial current}

\begin{itemize}
\item \tR{} \textbf{The framing.} ``Increasing the beam current is important, we
would just need to find the balance between a low (but not too low) dead time and
lowest possible carbon deposition.''
\basis{\prc{13}{5500446306}.}

\item \tC{} \textbf{Dose bookkeeping.} Counts track dose = current $\times$ live
time. For a fixed count target the dose is fixed --- current only sets how fast it
is delivered (real time = live time / $(1-\mathrm{DT})$) and where the electronics
operate. Refinement of the carbon coupling: carbon scales with dose per spot, so
faster delivery does \emph{not} add carbon --- the levers are cleanliness and fresh
spots, not lower current.
\basis{\prc{13}{5500449993}; \texttt{docs/eds\_parameter\_summary.md} \S5.}

\item \tC{} \textbf{The dead-time band --- and the open tension CALIBER must
settle.} The voltage-series plan operates at 25--30\% dead time (prediction P5 logs
the current achieving it per kV) --- the throughput-optimal band. The NIST
standards-based protocol (Newbury \& Ritchie) recommends $\le$10\%, stricter for
trace work, because coincidence/pile-up artifacts grow above $\approx$15\%. These
are in genuine tension: throughput says the Aug session used a tenth of the
detector's appetite ($\ge 10^4$ net Mg in $\approx$75~s instead of $\approx$12~min
at 10$\times$ the rate); the artifact argument says a nearly idle processor is the
cheapest insurance for a 0.5~wt\% element under a major peak. Standards themselves
need only \emph{in-range} dead time, not matched (the live-time clock corrects it).
\basis{Plan \S4 (\pr{13}); Newbury \& Ritchie [2] via the Edison best-practices
task
(\href{https://github.com/vertical-cloud-lab/caliber/blob/d6c29d1/outputs/eds-standards-best-practices/answer.md}{artifacts});
\texttt{docs/eds\_parameter\_summary.md} \S5.}

\item \tC{} \textbf{Throughput physics.} The SDD chain is paralyzable: stored rate
$=$ input rate $\times e^{-\,\mathrm{input} \cdot \tau}$, so beyond the optimum,
\emph{more} current stores \emph{fewer} counts; the amp time sets $\tau$. Dead time
is the busy fraction --- the observable that tells you where on the curve you sit.
\basis{\texttt{docs/eds\_daq.md} (\pr{6}); Amptek DPP theory [9].}

\item \tC{} \textbf{Measured current, not displayed current.} The set point
(``3.2~nA'') is a knob label --- it was never cup-verified. Faraday cup +
picoammeter routine: park in cup $\to$ read $\to$ acquire $\to$ read again; a close
pair proves stability and the average goes into the spectrum's probe-current field
(\texttt{\#PROBECUR} in \texttt{.msa} / DTSA-II properties). The k-ratio divides by
recorded dose, so \emph{known} current differences between standard and unknown
cancel exactly; \emph{unknown} ones become full-size wt\% errors. Warnings: display
vs cup off by tens of percent $\to$ aperture contamination / gun retune;
before/after differing $>$1\% $\to$ source drifted (use the mean; re-shoot
standards). Without a cup, a QC spectrum on a stable material (Si wafer) is the dose
proxy via count rate. \tR{} Mike confirmed the facility has both picoammeters and a
Faraday cup.
\basis{\prc{11}{5484533976} (routine; Ronnie's question \prc{11}{5484355320});
facility confirmation \prc{11}{5495597642}; Goldstein [1] ch.~16 / Newbury \&
Ritchie [2].}

\item \tR{} (with Xavier) \textbf{Search range to verify.} The outline's beam-current
range [20~pA, 10~mA] carries a ``?'' --- the instrument's real, cup-measured ladder
must bound the optimization space before any BO run uses it.
\basis{\texttt{CALIBER\_Outline.jpg}; open question 1 below.}
\end{itemize}

\subsection{Optimizing the current afterwards --- the open questions}

All six pre-staged in \pr{15} \tC{}, building on the framing above:

\begin{enumerate}[itemsep=2pt]
\item Map the Apreo's current ladder: cup-measured nA per preset, per menu kV.
\item Dead time vs current curve at 5~kV / 7.68~$\mu$s: which preset lands 25--30\%?
Which lands $\le$10\%? Verify the predicted $\approx$10$\times$ rate headroom over
the Aug session.
\item Pile-up onset: Al+Al sum-peak amplitude at 2.97~keV and Mg/Al net-ratio
stability as dead time climbs --- locate the artifact threshold for \emph{this}
detector and pick the CALIBER dead-time band from data (this simultaneously
adjudicates the P7 $\times$1.52 anomaly).
\item Stability: bracketing cup readings around each replicate --- drift per hour
and per session (a free QC drift chart).
\item The amp-time A/B (below), since amp time and current jointly set (dead time,
resolution).
\item Candidate CALIBER rule to validate: \emph{fix amp time at the
resolution-friendly setting; set current to the empirically chosen dead-time band;
measure it with the cup; record measured current $\times$ live time per spectrum.}
\end{enumerate}
\basis{\texttt{docs/eds\_parameter\_summary.md} \S5 (\pr{15}).}

\subsection{Amp (process) time --- partner knob, not a standalone parameter}

\begin{itemize}
\item \tR{} \textbf{The proposal.} After EDAX's ``under 4~eV'' resolution cost
between the slowest and fastest settings, Ronnie proposed treating amp time as
another parameter and testing 0.96 vs 7.68~$\mu$s to see whether slimming the Al and
Mg peaks helps noticeably.
\basis{\prc{13}{5593621690}; EDAX ladder 0.96/1.92/3.84/7.68~$\mu$s.}

\item \tC{} \textbf{The verdict --- worth one cheap A/B, not a parameter.} Direction
check first: \emph{longer} amp time = slimmer peaks, and 7.68~$\mu$s is already the
slim end --- 0.96~$\mu$s broadens, not slims. The ``4~eV'' is measured at Mn~K$\alpha$
(5.9~keV) and does not transfer: electronic noise is a constant-eV adder while the
Fano floor shrinks with photon energy, so Mg~K$\alpha$ broadens +10--25\%
(69 $\to$ 76--86~eV) and C~K$\alpha$ +20--47\%. Width was never the Al$\to$Mg lever
anyway --- the separation is 3.2 Al-FWHM, and the counts under Mg are continuum plus
the Al low-energy \emph{tail} (incomplete charge collection --- a peak-shape
property; short amp times risk ballistic deficit, making tails equal-or-worse). The
real trade is the count-rate ceiling: at 25--30\% dead time, $\approx$14~kcps stored
at 7.68~$\mu$s vs $\approx$81~kcps at 0.96~$\mu$s (5.7$\times$), with 1.92~$\mu$s a
good middle (+2~eV at Mg, 3.5$\times$ ceiling) --- but the ceiling only matters once
beam current rises to feed it. Hence: partner knob of probe current. The A/B runs as
an appendix block in the 5~kV session with pre-registered predictions A1--A4 (width
transfer; dose-invariance of net Mg; tail flat-or-worse; ceiling ratio
$\approx$5.7$\times$), kept out of the voltage series, which freezes amp time at
7.68~$\mu$s.
\basis{\prc{13}{5593623759}; EDAX dead-time tips sheet [8] (Fig.~1: Mn~K$\alpha$
124.1 $\to$ 128.0~eV); Fiori--Newbury FWHM scaling (Goldstein [1]).}
\end{itemize}

% =====================================================================
\section{Cross-cutting decisions the three parameters plug into}

\begin{itemize}
\item \tRC{} \textbf{Physical standards, purchased and vetted.} Order 13279
(approved 2026-09-08): high-purity Al evaporation pellet (Kurt J.\ Lesker; needs
polishing), Si chips, MgO~(100) wafer 10$\times$10$\times$0.5~mm (MTI). Claude
sourced candidate vendors and ran Edison vetting; Ronnie evaluated options, made the
selections, and purchased. MgO is why Mg gets a compound standard (pure Mg surfaces
poorly).
\basis{\prc{11}{5484853198}, \prc{11}{5485228187}, \prc{11}{5495597642};
\href{https://github.com/vertical-cloud-lab/caliber/blob/d6c29d1/outputs/eds-standards-material-vetting/answer.md}{Edison vetting artifacts}.}

\item \tRC{} \textbf{Standard reuse rules.} Ronnie asked whether one standard with
given settings serves all future samples. Answer: a standard is defined by
(element/line, kV, detector, process time) and is reusable \emph{indefinitely} given
per-session QC --- energy calibration, FWHM, counts per nA$\cdot$s on a QC material,
Duane--Hunt endpoint. Acquisition parameters split into: \emph{identical} (detector,
kV set point, process time, geometry/working distance, coating scheme);
\emph{measured every acquisition} (probe current, live time --- the dose); and
\emph{in-range, not matched} (dead time, count rate).
\basis{\prc{11}{5481700165} (question), \prc{11}{5481703178} (rules);
Newbury \& Ritchie [2]; Goldstein [1] ch.~16.}

\item \tMS{} \textbf{Geometry discipline: ``focusing to Z height.''} Acquire
standards at a specific stage height, then acquire every unknown at exactly the same
height (item to walk through with Mike). Physics behind it \tC{}: working distance
sets the take-off angle, and the absorption path scales as $1/\sin\psi$ --- the
matrix correction is calibrated for one geometry. \tR{} Corollary: keep everything at
the 10~mm eucentric height --- which rules out dual EDS/EBSD sessions (EBSD needs
15~mm and a $70^\circ$ tilt).
\basis{Mike's notes \issc{5513720218}; Ronnie's eucentric-height reasoning
\prc{11}{5441452199}.}

\item \tMS{} \textbf{Type standards from characterized samples.} Once a sample is
well-characterized against the pure-element standards or by ICP-MS/XRF, use it as
the working standard for other samples of the same alloy --- cheaper sessions, same
traceability.
\basis{Mike's notes \issc{5513720218}.}

\item \tMS{} \textbf{Honest limits.} EDS has a floor; Mike agreed the standards +
matched-conditions path is the best available route to high-accuracy wt\% for
sub-1-wt\% constituents. Where a number really matters, the ICP referee decides.
\basis{Mike's notes \issc{5513720218}.}

\item \tRC{} \textbf{Software chain.} DTSA-II for wt\% (reads \texttt{.msa}
natively, $\varphi(\rho z)$ matrix correction, bundle/quant wizards) --- adopted by
Ronnie after Claude's assessment that eXSpy has no bulk-SEM matrix correction (its
quantification is TEM thin-film only) and that CalcZAF consumes k-ratios, not
spectra. DTSA-II also runs programmatically/headless (verified in CI), so the
pipeline can be automated. To do \tR{}: identify the facility's EDAX detector model
and import its detector definition into DTSA-II.
\basis{\issc{5359231847}; Edison tools comparison
(\href{https://github.com/vertical-cloud-lab/caliber/blob/d6c29d1/outputs/edison_eds_tools_comparison/answer.md}{artifacts});
headless proof \issc{5629977539}; detector-import note \issc{5534375327}.}

\item \tRC{} \textbf{Charging and coating for the MgO standard.} MgO is an insulator
and needs a carbon coat; the BYU coater is in storage until end of October. VPSEM is
\emph{not} an acceptable quantitative workaround (the beam skirt scatters a
pressure-dependent fraction of current hundreds of $\mu$m off-probe) --- Claude/Edison
concluded it, and Ronnie independently double-checked the literature and concurred.
Ronnie's exploratory idea on charging mitigation: a grounded conductive layer placed
below the X-ray production depth but within electron range. Per-spectrum charging
check either way: the Duane--Hunt endpoint (P6).
\basis{\prc{11}{5515514463}, \prc{11}{5534844556};
\href{https://github.com/vertical-cloud-lab/caliber/blob/d6c29d1/outputs/vpsem-low-vacuum-eds/answer.md}{Edison VPSEM artifacts},
\href{https://github.com/vertical-cloud-lab/caliber/blob/d6c29d1/outputs/insulator-charging-depth-grounding/answer.md}{charging-depth artifacts}.}
\end{itemize}

% =====================================================================
\section*{References}
\begin{enumerate}[label={[\arabic*]}, itemsep=2pt]
\item Goldstein, Newbury, Michael, Ritchie, Scott \& Joy, \emph{Scanning Electron
Microscopy and X-Ray Microanalysis}, 4th ed., Springer (2018). Interaction volume
\S3.3.3.1; QC ch.~16; k-ratio procedure ch.~19--20; the SDD count-target passage
($>10^6$/$10^7$/$10^8$) quoted by Ronnie in \prc{13}{5574391518}.
\item Newbury \& Ritchie, ``Performing elemental microanalysis with high accuracy
and high precision by SEM/SDD-EDS,'' \emph{J.\ Mater.\ Sci.} 50:493--518 (2015).
\href{https://doi.org/10.1007/s10853-014-8685-2}{doi:10.1007/s10853-014-8685-2}.
The program's core protocol paper.
\item Statham, ``Limitations to Accuracy in Extracting Characteristic Line
Intensities From X-Ray Spectra,'' \emph{J.\ Res.\ NIST} 107:531 (2002).
\href{https://pmc.ncbi.nlm.nih.gov/articles/PMC4863855/}{open access}.
\item Currie, ``Limits for qualitative detection and quantitative determination,''
\emph{Anal.\ Chem.} 40:586 (1968).
\href{https://doi.org/10.1021/ac60259a007}{doi:10.1021/ac60259a007}.
\item Ziebold, ``Precision and sensitivity in electron microprobe analysis,''
\emph{Anal.\ Chem.} 39:858 (1967).
\href{https://doi.org/10.1021/ac60252a028}{doi:10.1021/ac60252a028}.
\item Kanaya \& Okayama, ``Penetration and energy-loss theory of electrons in solid
targets,'' \emph{J.\ Phys.\ D} 5:43 (1972). Used with the production-depth
refinement $E_0^{1.67} \to (E_0^{1.67} - E_c^{1.67})$.
\item Philibert absorption model / $f(\chi)$, with Elam mass-absorption
coefficients via \texttt{xraydb} --- as implemented in
\texttt{scripts/eds\_voltage\_predictions.py} (\pr{13}).
\item EDAX, ``Dead Time Optimization for High Productivity Data Analysis''
(tips sheet) --- amp-time ladder, measured resolution/throughput trade.
\item Amptek application notes: ``Digital Pulse Processors --- Theory of
Operation'' and the SDD note --- basis of \texttt{docs/eds\_daq.md} (\pr{6}).
\item Ritchie, NIST DTSA-II,
\href{https://www.cstl.nist.gov/div837/837.02/epq/dtsa2/index.html}{cstl.nist.gov/div837/837.02/epq/dtsa2}.
\item Newbury, Swyt \& Myklebust, ``Standardless quantitative electron probe
microanalysis with EDS: is it worth the risk?'' \emph{Anal.\ Chem.} 67:1866 (1995).
\href{https://doi.org/10.1021/ac00107a017}{doi:10.1021/ac00107a017}.
\item Edison Scientific literature tasks committed under
\href{https://github.com/vertical-cloud-lab/caliber/tree/claude/issue-1-20260909-0027/outputs}{\texttt{outputs/}}
on the \pr{15} branch: parameter derivations, standards protocol and best
practices, material vetting, reading list, insulator charging, VPSEM, tools
comparison, Mg evaporation.
\end{enumerate}

\end{document}
